\documentclass[11pt,a4paper]{article}

% ---- Packages ----
\usepackage[utf8]{inputenc}
\usepackage[T1]{fontenc}
\usepackage{amsmath,amssymb,amsthm}
\usepackage{mathtools}
\usepackage[margin=1in,headheight=14pt]{geometry}
\usepackage{hyperref}
\usepackage{enumitem}
\usepackage{xcolor}
\usepackage{tcolorbox}
\usepackage{fancyhdr}
\usepackage{booktabs}

% ---- Hyperref ----
\hypersetup{
    colorlinks=true,
    linkcolor=red,
    citecolor=blue,
    urlcolor=blue
}

% ---- Header ----
\pagestyle{fancy}
\fancyhf{}
\fancyhead[L]{\textit{Real Analysis Review}}
\fancyhead[R]{\thepage}
\renewcommand{\headrulewidth}{0.4pt}

% ---- Theorem Environments ----
\theoremstyle{definition}
\newtheorem{definition}{Definition}[section]
\newtheorem{example}{Example}[section]
\newtheorem{exercise}{Exercise}[section]

\theoremstyle{plain}
\newtheorem{theorem}{Theorem}[section]
\newtheorem{lemma}[theorem]{Lemma}
\newtheorem{proposition}[theorem]{Proposition}
\newtheorem{corollary}[theorem]{Corollary}

\theoremstyle{remark}
\newtheorem{remark}{Remark}[section]

% ---- Colored Boxes ----
\tcbuselibrary{theorems,skins,breakable}

\newtcolorbox{keyresult}[1][]{
    colback=green!5!white,
    colframe=green!50!black,
    fonttitle=\bfseries,
    title=#1,
    breakable
}

\newtcolorbox{rlconnection}[1][]{
    colback=blue!5!white,
    colframe=blue!50!black,
    fonttitle=\bfseries,
    title=#1,
    breakable
}

\newtcolorbox{warning}[1][]{
    colback=red!5!white,
    colframe=red!50!black,
    fonttitle=\bfseries,
    title=#1,
    breakable
}

% ---- Operators ----
\newcommand{\R}{\mathbb{R}}
\newcommand{\N}{\mathbb{N}}
\newcommand{\Q}{\mathbb{Q}}
\newcommand{\Z}{\mathbb{Z}}

% ---- Title ----
\title{\textbf{Real Analysis Review}\\[10pt]
\large Sequences, Continuity, and Convergence\\[5pt]
\normalsize Phase 0: Prerequisites Review $\mid$ Document 1 of 3\\
\normalsize Primary texts: Rudin, \emph{Principles of Mathematical Analysis}; Abbott, \emph{Understanding Analysis}}
\author{Curriculum: A Rigorous Learning Path to Multi-Agent Deep RL}
\date{February 2026}

\begin{document}

\maketitle

\begin{abstract}
This document is a self-contained refresher of undergraduate real analysis, aimed at
readers who have seen this material before but need the key results recalled with full
formal statements and complete proofs. The selection of topics is driven entirely by
what is needed for measure theory, probability, and functional analysis in Phase~01
and beyond: the completeness of~$\R$, convergence of sequences and series, topology
of the real line, continuity, differentiation, Riemann integration, sequences of
functions, and metric spaces. Every theorem is proved. The reader who works through
this review will be fully prepared for the construction of $\sigma$-algebras, measures,
and the Lebesgue integral.
\end{abstract}

\tableofcontents
\newpage

%% ============================================================
%% SECTION 1: INTRODUCTION
%% ============================================================
\section{Introduction}
\label{sec:introduction}

This document is not a first course in real analysis. It is a focused refresher for
readers who have taken undergraduate analysis but find the details ``dusty.'' Our goal
is to recall --- with full proofs --- exactly the results that will be needed when we
move to measure theory and beyond.

Here is what we cover and why each topic matters:

\begin{enumerate}[nosep]
    \item \textbf{The real number system} (Section~\ref{sec:real_numbers}). The
    completeness axiom (least upper bound property) is the single most important
    property of~$\R$. It drives the construction of Lebesgue measure, ensures the
    existence of suprema needed to define outer measures, and underlies every
    convergence theorem in the curriculum.

    \item \textbf{Sequences and series} (Section~\ref{sec:sequences_series}). The
    monotone convergence theorem for sequences, Bolzano--Weierstrass, Cauchy
    completeness, and $\limsup/\liminf$ are the workhorses of measure-theoretic
    arguments. The $\limsup$ of a sequence of sets, for instance, is built directly
    from the $\limsup$ concept reviewed here.

    \item \textbf{Topology of the real line} (Section~\ref{sec:topology}). Open sets,
    closed sets, and compactness are the language in which $\sigma$-algebras and Borel
    sets are defined. The Heine--Borel theorem characterizes compact subsets of~$\R^n$.

    \item \textbf{Continuity} (Section~\ref{sec:continuity}). Continuous functions
    on compact sets attain their bounds and are uniformly continuous --- facts used
    repeatedly in measure theory and functional analysis.

    \item \textbf{Differentiation} (Section~\ref{sec:differentiation}). The mean value
    theorem and Taylor's theorem are essential for optimization theory and appear in
    many $\varepsilon$-$\delta$ arguments throughout the curriculum.

    \item \textbf{Riemann integration} (Section~\ref{sec:riemann}). Understanding the
    Riemann integral --- and its limitations --- motivates the passage to the Lebesgue
    integral in Phase~01.

    \item \textbf{Sequences and series of functions}
    (Section~\ref{sec:function_sequences}). Uniform convergence and its interchange
    theorems are the starting point; the dominated and monotone convergence theorems
    of measure theory vastly generalize these results.

    \item \textbf{Metric spaces} (Section~\ref{sec:metric_spaces}). The abstract
    setting of metric spaces is where functional analysis begins. The contraction
    mapping theorem proved here is directly used to establish convergence of value
    iteration in reinforcement learning.
\end{enumerate}

\begin{rlconnection}[From Real Analysis to Measure Theory]
Measure theory generalizes Riemann integration to handle limits of functions that
Riemann cannot. The concepts in this review --- $\sup/\inf$, sequences, limits,
open/closed sets, compactness --- are the raw material from which $\sigma$-algebras,
measures, and Lebesgue integration are built.
\end{rlconnection}

%% ============================================================
%% SECTION 2: THE REAL NUMBER SYSTEM
%% ============================================================
\section{The Real Number System}
\label{sec:real_numbers}

\subsection{Ordered Field Axioms}
\label{subsec:ordered_field}

We take the real numbers $\R$ as given, satisfying the following axioms.

\begin{definition}[Ordered Field]
\label{def:ordered_field}
The real numbers $\R$ form an \emph{ordered field}: a set equipped with
operations $+$ and $\cdot$ and a total order $\le$ satisfying:
\begin{enumerate}[nosep,label=(\roman*)]
    \item \textbf{Field axioms.} $(\R, +, \cdot)$ is a field: addition and
    multiplication are commutative, associative, and distributive, with identity
    elements $0$ and $1$ respectively, additive inverses for all elements, and
    multiplicative inverses for all nonzero elements.
    \item \textbf{Order axioms.} $\le$ is a total order compatible with the field
    operations: if $a \le b$ then $a + c \le b + c$ for all $c$; if $0 \le a$ and
    $0 \le b$ then $0 \le ab$.
\end{enumerate}
\end{definition}

\begin{remark}
\label{rem:rational_field}
The rationals $\Q$ also form an ordered field. What distinguishes $\R$ from $\Q$
is the completeness axiom below.
\end{remark}

\subsection{The Completeness Axiom}
\label{subsec:completeness}

\begin{definition}[Upper Bound, Supremum]
\label{def:supremum}
Let $S \subseteq \R$ be nonempty.
\begin{enumerate}[nosep,label=(\roman*)]
    \item An element $u \in \R$ is an \emph{upper bound} for $S$ if $s \le u$
    for every $s \in S$. We say $S$ is \emph{bounded above} if such a $u$ exists.
    \item An element $\alpha \in \R$ is the \emph{supremum} (or \emph{least upper
    bound}) of $S$, written $\alpha = \sup S$, if:
    \begin{enumerate}[nosep,label=(\alph*)]
        \item $\alpha$ is an upper bound for $S$, and
        \item if $u$ is any upper bound for $S$, then $\alpha \le u$.
    \end{enumerate}
\end{enumerate}
\end{definition}

\begin{definition}[Lower Bound, Infimum]
\label{def:infimum}
Dually, $\ell \in \R$ is a \emph{lower bound} for $S$ if $\ell \le s$ for every
$s \in S$. The \emph{infimum} $\inf S$ is the greatest lower bound.
\end{definition}

\begin{remark}
\label{rem:sup_characterization}
The supremum $\alpha = \sup S$ is characterized by: $\alpha$ is an upper bound
for $S$, and for every $\varepsilon > 0$ there exists $s \in S$ with
$s > \alpha - \varepsilon$. This characterization is used constantly.
\end{remark}

\begin{theorem}[Completeness Axiom / Least Upper Bound Property]
\label{thm:completeness}
Every nonempty subset of\/ $\R$ that is bounded above has a supremum in~$\R$.
\end{theorem}

This is an \emph{axiom}: it is not proved from more basic principles, but rather
is the defining property that distinguishes $\R$ from $\Q$.

\begin{corollary}[Greatest Lower Bound Property]
\label{cor:glb}
Every nonempty subset of\/ $\R$ that is bounded below has an infimum in~$\R$.
\end{corollary}

\begin{proof}
Let $S \subseteq \R$ be nonempty and bounded below. Define
$-S = \{-s : s \in S\}$. Then $-S$ is nonempty and bounded above
(if $\ell$ is a lower bound for $S$, then $-\ell$ is an upper bound for $-S$).
By the completeness axiom, $\alpha = \sup(-S)$ exists. We claim $\inf S = -\alpha$.

Since $\alpha$ is an upper bound for $-S$, we have $-s \le \alpha$ for all $s \in S$,
hence $s \ge -\alpha$ for all $s \in S$, so $-\alpha$ is a lower bound for $S$.
If $\ell$ is any lower bound for $S$, then $-\ell$ is an upper bound for $-S$,
so $\alpha \le -\ell$, hence $\ell \le -\alpha$. Thus $-\alpha$ is the greatest
lower bound.
\end{proof}

\subsection{Archimedean Property and Density of Rationals}
\label{subsec:archimedean}

\begin{theorem}[Archimedean Property]
\label{thm:archimedean}
For every $x \in \R$, there exists $n \in \N$ with $n > x$.
Equivalently, for every $\varepsilon > 0$ and every $M > 0$, there exists
$n \in \N$ with $n\varepsilon > M$.
\end{theorem}

\begin{proof}
Suppose for contradiction that $\N$ is bounded above in $\R$. By the completeness
axiom, $\alpha = \sup \N$ exists. Then $\alpha - 1$ is not an upper bound for $\N$
(since $\alpha$ is the \emph{least} upper bound), so there exists $m \in \N$ with
$m > \alpha - 1$. But then $m + 1 > \alpha$ with $m + 1 \in \N$, contradicting the
fact that $\alpha$ is an upper bound for $\N$.

The second formulation follows by applying the first to $x = M/\varepsilon$.
\end{proof}

\begin{theorem}[Density of Rationals]
\label{thm:density_Q}
For every $a, b \in \R$ with $a < b$, there exists $q \in \Q$ with $a < q < b$.
\end{theorem}

\begin{proof}
By the Archimedean property, choose $n \in \N$ with $n(b - a) > 1$, i.e.,
$nb - na > 1$. Again by the Archimedean property, the sets
$\{k \in \Z : k > na\}$ and $\{k \in \Z : k < na\}$ are both nonempty, so we may
define $m = \min\{k \in \Z : k > na\}$ (this minimum exists because the positive
integers exceeding $na$ form a nonempty subset of $\N$, which has a least element
by the well-ordering principle). Then $m > na$ and $m - 1 \le na$, so
$m \le na + 1 < nb$ (since $nb - na > 1$). Setting $q = m/n$ gives
$a < q < b$ with $q \in \Q$.
\end{proof}

\begin{theorem}[Nested Intervals Property]
\label{thm:nested_intervals}
Let $I_n = [a_n, b_n]$ be a sequence of closed bounded intervals with
$I_1 \supseteq I_2 \supseteq I_3 \supseteq \cdots$. Then
$\bigcap_{n=1}^\infty I_n \neq \emptyset$.
\end{theorem}

\begin{proof}
The nesting condition gives $a_1 \le a_2 \le \cdots$ and
$b_1 \ge b_2 \ge \cdots$ with $a_n \le b_n$ for all $n$. In particular,
$a_n \le b_1$ for all $n$, so the set $\{a_n : n \in \N\}$ is bounded above.
By the completeness axiom, $\alpha = \sup\{a_n : n \in \N\}$ exists.

For every $n$, we have $a_n \le \alpha$ (since $\alpha$ is an upper bound) and
$\alpha \le b_n$ (since $b_n$ is an upper bound for $\{a_k\}$: for $k \le n$,
$a_k \le a_n \le b_n$; for $k > n$, $a_k \le b_k \le b_n$). Therefore
$\alpha \in [a_n, b_n]$ for every $n$, giving
$\alpha \in \bigcap_{n=1}^\infty I_n$.
\end{proof}

\begin{remark}
\label{rem:nested_open}
The conclusion fails for open intervals: $\bigcap_{n=1}^\infty (0, 1/n) = \emptyset$.
It also fails for unbounded closed sets: $\bigcap_{n=1}^\infty [n, \infty) = \emptyset$.
Both closure and boundedness are essential.
\end{remark}

%% ============================================================
%% SECTION 3: SEQUENCES AND SERIES
%% ============================================================
\section{Sequences and Series}
\label{sec:sequences_series}

\subsection{Convergence of Sequences}
\label{subsec:seq_convergence}

\begin{definition}[Convergence]
\label{def:convergence}
A sequence $(a_n)_{n=1}^\infty$ in $\R$ \emph{converges} to $L \in \R$, written
$\lim_{n \to \infty} a_n = L$ or $a_n \to L$, if for every $\varepsilon > 0$
there exists $N \in \N$ such that $|a_n - L| < \varepsilon$ for all $n \ge N$.
A sequence that does not converge is \emph{divergent}.
\end{definition}

\begin{theorem}[Uniqueness of Limits]
\label{thm:unique_limit}
If $a_n \to L$ and $a_n \to M$, then $L = M$.
\end{theorem}

\begin{proof}
Let $\varepsilon > 0$. There exist $N_1, N_2 \in \N$ such that
$|a_n - L| < \varepsilon/2$ for $n \ge N_1$ and $|a_n - M| < \varepsilon/2$
for $n \ge N_2$. For $n \ge \max(N_1, N_2)$:
\[
|L - M| \le |L - a_n| + |a_n - M| < \frac{\varepsilon}{2} + \frac{\varepsilon}{2}
= \varepsilon.
\]
Since $|L - M| < \varepsilon$ for every $\varepsilon > 0$, we conclude $L = M$.
\end{proof}

\begin{theorem}[Algebraic Limit Theorems]
\label{thm:alg_limit}
Suppose $a_n \to a$ and $b_n \to b$. Then:
\begin{enumerate}[nosep,label=(\roman*)]
    \item $a_n + b_n \to a + b$.
    \item $a_n \cdot b_n \to a \cdot b$.
    \item If $b \neq 0$, then $a_n / b_n \to a / b$ (for $n$ large enough that
    $b_n \neq 0$).
\end{enumerate}
\end{theorem}

\begin{proof}
We prove~(ii); the proofs of~(i) and~(iii) are similar.

Since $(a_n)$ converges, it is bounded: there exists $M > 0$ with $|a_n| \le M$
for all $n$. Let $\varepsilon > 0$. Choose $N_1$ so that
$|a_n - a| < \varepsilon/(2(|b| + 1))$ for $n \ge N_1$, and $N_2$ so that
$|b_n - b| < \varepsilon/(2M)$ for $n \ge N_2$. For $n \ge \max(N_1, N_2)$:
\begin{align*}
|a_n b_n - ab|
&= |a_n b_n - a_n b + a_n b - ab| \\
&\le |a_n||b_n - b| + |b||a_n - a| \\
&\le M \cdot \frac{\varepsilon}{2M} + |b| \cdot \frac{\varepsilon}{2(|b|+1)} \\
&< \frac{\varepsilon}{2} + \frac{\varepsilon}{2} = \varepsilon. \qedhere
\end{align*}
\end{proof}

\subsection{Monotone Convergence and Bolzano--Weierstrass}
\label{subsec:monotone_bw}

\begin{theorem}[Monotone Convergence Theorem for Sequences]
\label{thm:mct_seq}
Every bounded monotone sequence converges. Specifically:
\begin{enumerate}[nosep,label=(\roman*)]
    \item If $(a_n)$ is increasing and bounded above, then
    $a_n \to \sup\{a_n : n \in \N\}$.
    \item If $(a_n)$ is decreasing and bounded below, then
    $a_n \to \inf\{a_n : n \in \N\}$.
\end{enumerate}
\end{theorem}

\begin{proof}
We prove~(i). Let $\alpha = \sup\{a_n : n \in \N\}$, which exists by the
completeness axiom. Let $\varepsilon > 0$. Since $\alpha - \varepsilon$ is not
an upper bound for $\{a_n\}$, there exists $N$ with $a_N > \alpha - \varepsilon$.
Since $(a_n)$ is increasing, for $n \ge N$:
\[
\alpha - \varepsilon < a_N \le a_n \le \alpha,
\]
so $|a_n - \alpha| < \varepsilon$. Part~(ii) follows by applying~(i) to $(-a_n)$.
\end{proof}

\begin{theorem}[Bolzano--Weierstrass]
\label{thm:bolzano_weierstrass}
Every bounded sequence in $\R$ has a convergent subsequence.
\end{theorem}

\begin{proof}
Let $(a_n)$ be bounded. We construct a monotone subsequence, which converges by
Theorem~\ref{thm:mct_seq}.

Call index $n$ a \emph{peak} if $a_n \ge a_m$ for all $m > n$. There are two cases:

\textbf{Case 1:} There are infinitely many peaks $n_1 < n_2 < n_3 < \cdots$.
Then $a_{n_1} \ge a_{n_2} \ge a_{n_3} \ge \cdots$ is a decreasing subsequence.
It is bounded (since $(a_n)$ is bounded), hence convergent by
Theorem~\ref{thm:mct_seq}.

\textbf{Case 2:} There are only finitely many peaks. Let $N$ be larger than
all peak indices. Then no index $n \ge N$ is a peak, meaning for every $n \ge N$
there exists $m > n$ with $a_m > a_n$. Starting with $n_1 = N$, inductively choose
$n_{k+1} > n_k$ with $a_{n_{k+1}} > a_{n_k}$. This gives a strictly increasing
subsequence. It is bounded above (since $(a_n)$ is bounded), hence convergent by
Theorem~\ref{thm:mct_seq}.
\end{proof}

\subsection{Cauchy Sequences and Completeness}
\label{subsec:cauchy}

\begin{definition}[Cauchy Sequence]
\label{def:cauchy}
A sequence $(a_n)$ in $\R$ is \emph{Cauchy} if for every $\varepsilon > 0$ there
exists $N \in \N$ such that $|a_n - a_m| < \varepsilon$ for all $n, m \ge N$.
\end{definition}

\begin{theorem}[Completeness of $\R$]
\label{thm:cauchy_complete}
A sequence in $\R$ converges if and only if it is Cauchy.
\end{theorem}

\begin{proof}
$(\Rightarrow)$ Suppose $a_n \to L$. Given $\varepsilon > 0$, choose $N$ so that
$|a_n - L| < \varepsilon/2$ for $n \ge N$. For $n, m \ge N$:
\[
|a_n - a_m| \le |a_n - L| + |L - a_m| < \varepsilon.
\]

$(\Leftarrow)$ Suppose $(a_n)$ is Cauchy. First, $(a_n)$ is bounded: choose $N$
so that $|a_n - a_m| < 1$ for $n, m \ge N$; then $|a_n| \le \max(|a_1|, \ldots,
|a_{N-1}|, |a_N| + 1)$ for all $n$.

By Bolzano--Weierstrass (Theorem~\ref{thm:bolzano_weierstrass}), $(a_n)$ has a
convergent subsequence $a_{n_k} \to L$. We show $a_n \to L$. Given
$\varepsilon > 0$, choose $N_1$ so that $|a_n - a_m| < \varepsilon/2$ for
$n, m \ge N_1$, and choose $K$ so that $|a_{n_k} - L| < \varepsilon/2$ for
$k \ge K$. Pick $k \ge K$ with $n_k \ge N_1$. Then for $n \ge N_1$:
\[
|a_n - L| \le |a_n - a_{n_k}| + |a_{n_k} - L|
< \frac{\varepsilon}{2} + \frac{\varepsilon}{2} = \varepsilon. \qedhere
\]
\end{proof}

\subsection{Limit Superior and Limit Inferior}
\label{subsec:limsup_liminf}

\begin{definition}[$\limsup$ and $\liminf$]
\label{def:limsup_liminf}
For a bounded sequence $(a_n)$, define:
\[
\limsup_{n \to \infty} a_n = \lim_{n \to \infty}
\Bigl(\sup_{k \ge n} a_k\Bigr), \qquad
\liminf_{n \to \infty} a_n = \lim_{n \to \infty}
\Bigl(\inf_{k \ge n} a_k\Bigr).
\]
The sequences $\sup_{k \ge n} a_k$ (decreasing) and $\inf_{k \ge n} a_k$
(increasing) are monotone and bounded, so both limits exist by
Theorem~\ref{thm:mct_seq}.
\end{definition}

\begin{proposition}[Properties of $\limsup$ and $\liminf$]
\label{prop:limsup_props}
Let $(a_n)$ be a bounded sequence. Then:
\begin{enumerate}[nosep,label=(\roman*)]
    \item $\liminf_{n\to\infty} a_n \le \limsup_{n\to\infty} a_n$.
    \item $(a_n)$ converges if and only if
    $\liminf_{n\to\infty} a_n = \limsup_{n\to\infty} a_n$, and in that case
    the limit equals this common value.
    \item $\limsup_{n\to\infty} a_n = L$ if and only if: for every
    $\varepsilon > 0$, (a)~$a_n < L + \varepsilon$ for all sufficiently
    large $n$, and (b)~$a_n > L - \varepsilon$ for infinitely many $n$.
\end{enumerate}
\end{proposition}

\begin{proof}
(i) For each $n$, $\inf_{k \ge n} a_k \le \sup_{k \ge n} a_k$. Taking limits
preserves the inequality.

(ii) $(\Rightarrow)$ Suppose $a_n \to L$. Given $\varepsilon > 0$, for large $n$
we have $L - \varepsilon < a_k < L + \varepsilon$ for all $k \ge n$, so
$L - \varepsilon \le \inf_{k \ge n} a_k \le \sup_{k \ge n} a_k \le L + \varepsilon$.
Taking the limit gives $L - \varepsilon \le \liminf a_n \le \limsup a_n
\le L + \varepsilon$ for all $\varepsilon > 0$, so both equal $L$.

$(\Leftarrow)$ Suppose $\liminf a_n = \limsup a_n = L$. For any $\varepsilon > 0$,
choose $N$ so that $|\sup_{k \ge n} a_k - L| < \varepsilon$ and
$|\inf_{k \ge n} a_k - L| < \varepsilon$ for $n \ge N$. Then for $n \ge N$:
\[
L - \varepsilon < \inf_{k \ge n} a_k \le a_n \le \sup_{k \ge n} a_k < L + \varepsilon,
\]
so $|a_n - L| < \varepsilon$.

(iii) Let $s_n = \sup_{k \ge n} a_k$ and $L = \lim s_n$. For~(a): given
$\varepsilon > 0$, choose $N$ with $s_N < L + \varepsilon$. Then $a_n \le s_N <
L + \varepsilon$ for all $n \ge N$. For~(b): given $\varepsilon > 0$, for every $n$
we have $s_n \ge L$ (since $s_n$ is decreasing to $L$), so there exists $k \ge n$
with $a_k > s_n - \varepsilon \ge L - \varepsilon$. Since $n$ is arbitrary, this
gives infinitely many such $k$.

Conversely, suppose (a) and (b) hold for some value $L'$. From~(a), $s_n \le
L' + \varepsilon$ for large $n$, so $\limsup a_n \le L' + \varepsilon$.
From~(b), for each $n$ there exists $k \ge n$ with $a_k > L' - \varepsilon$,
so $s_n \ge L' - \varepsilon$, giving $\limsup a_n \ge L' - \varepsilon$.
Since $\varepsilon$ is arbitrary, $\limsup a_n = L'$.
\end{proof}

\subsection{Series}
\label{subsec:series}

\begin{definition}[Convergence of Series]
\label{def:series}
Let $(a_n)_{n=1}^\infty$ be a sequence. The series $\sum_{n=1}^\infty a_n$
\emph{converges} if the sequence of partial sums $S_N = \sum_{n=1}^N a_n$
converges. The series \emph{converges absolutely} if $\sum_{n=1}^\infty |a_n|$
converges.
\end{definition}

\begin{theorem}[Comparison Test]
\label{thm:comparison}
If $0 \le a_n \le b_n$ for all $n$ and $\sum b_n$ converges, then $\sum a_n$
converges.
\end{theorem}

\begin{proof}
The partial sums $S_N = \sum_{n=1}^N a_n$ form an increasing sequence (since
$a_n \ge 0$) bounded above by $\sum_{n=1}^\infty b_n$. By the monotone
convergence theorem (Theorem~\ref{thm:mct_seq}), $(S_N)$ converges.
\end{proof}

\begin{theorem}[Absolute Convergence Implies Convergence]
\label{thm:abs_conv}
If $\sum_{n=1}^\infty |a_n|$ converges, then $\sum_{n=1}^\infty a_n$ converges.
\end{theorem}

\begin{proof}
Absolute convergence means the partial sums of $|a_n|$ form a Cauchy sequence.
For $m > n$:
\[
\Bigl|\sum_{k=n+1}^{m} a_k\Bigr| \le \sum_{k=n+1}^{m} |a_k|,
\]
so the partial sums of $a_n$ are also Cauchy, hence convergent by
Theorem~\ref{thm:cauchy_complete}.
\end{proof}

\begin{theorem}[Ratio Test]
\label{thm:ratio}
Let $(a_n)$ be a sequence with $a_n \neq 0$ for all $n$. If
$\limsup_{n\to\infty} |a_{n+1}/a_n| < 1$, then $\sum a_n$ converges absolutely.
If $\liminf_{n\to\infty} |a_{n+1}/a_n| > 1$, then $\sum a_n$ diverges.
\end{theorem}

\begin{proof}
Suppose $\limsup |a_{n+1}/a_n| = L < 1$. Choose $r$ with $L < r < 1$.
By Proposition~\ref{prop:limsup_props}(iii), there exists $N$ with
$|a_{n+1}/a_n| < r$ for all $n \ge N$. Then $|a_n| \le |a_N| r^{n-N}$ for
$n \ge N$, and $\sum |a_N| r^{n-N}$ converges (geometric series with ratio
$r < 1$). By the comparison test, $\sum |a_n|$ converges.

If $\liminf |a_{n+1}/a_n| > 1$, then $|a_{n+1}| > |a_n|$ for large $n$, so
$a_n \not\to 0$ and the series diverges.
\end{proof}

\begin{theorem}[Root Test]
\label{thm:root}
Let $\alpha = \limsup_{n\to\infty} |a_n|^{1/n}$. If $\alpha < 1$, then
$\sum a_n$ converges absolutely. If $\alpha > 1$, then $\sum a_n$ diverges.
\end{theorem}

\begin{proof}
If $\alpha < 1$, choose $r$ with $\alpha < r < 1$. By
Proposition~\ref{prop:limsup_props}(iii), $|a_n|^{1/n} < r$ for all
sufficiently large $n$, so $|a_n| < r^n$. The comparison test with the geometric
series $\sum r^n$ gives absolute convergence.

If $\alpha > 1$, then $|a_n|^{1/n} > 1$ for infinitely many $n$
(Proposition~\ref{prop:limsup_props}(iii)), so $|a_n| > 1$ infinitely often,
meaning $a_n \not\to 0$ and the series diverges.
\end{proof}

\begin{remark}[Conditional Convergence]
\label{rem:conditional}
A series that converges but does not converge absolutely is called
\emph{conditionally convergent}. The classic example is the alternating harmonic
series $\sum (-1)^{n+1}/n$. By the Riemann rearrangement theorem, any
conditionally convergent series can be rearranged to converge to any prescribed
value (or to diverge). This pathology motivates the focus on absolute convergence
and non-negative series in measure theory.
\end{remark}

%% ============================================================
%% SECTION 4: TOPOLOGY OF THE REAL LINE
%% ============================================================
\section{Topology of the Real Line}
\label{sec:topology}

\subsection{Open and Closed Sets}
\label{subsec:open_closed}

\begin{definition}[Open and Closed Sets]
\label{def:open_closed}
A set $U \subseteq \R$ is \emph{open} if for every $x \in U$ there exists
$\varepsilon > 0$ such that $(x - \varepsilon, x + \varepsilon) \subseteq U$.
A set $F \subseteq \R$ is \emph{closed} if its complement $\R \setminus F$ is
open.
\end{definition}

\begin{proposition}
\label{prop:open_closed_basics}
\begin{enumerate}[nosep,label=(\roman*)]
    \item The empty set $\emptyset$ and $\R$ are both open and closed.
    \item An arbitrary union of open sets is open.
    \item A finite intersection of open sets is open.
    \item An arbitrary intersection of closed sets is closed.
    \item A finite union of closed sets is closed.
\end{enumerate}
\end{proposition}

\begin{proof}
(i) $\emptyset$ is vacuously open; $\R$ is open since every $\varepsilon$-neighborhood
is contained in $\R$. Their complements give the closed assertions.

(ii) If $x \in \bigcup_\alpha U_\alpha$ with each $U_\alpha$ open, then $x \in U_\beta$
for some $\beta$, so some $\varepsilon$-neighborhood of $x$ is contained in $U_\beta
\subseteq \bigcup_\alpha U_\alpha$.

(iii) If $x \in \bigcap_{i=1}^n U_i$ with each $U_i$ open, choose $\varepsilon_i > 0$
with $(x - \varepsilon_i, x + \varepsilon_i) \subseteq U_i$. Then
$\varepsilon = \min(\varepsilon_1, \ldots, \varepsilon_n) > 0$ and
$(x - \varepsilon, x + \varepsilon) \subseteq \bigcap_{i=1}^n U_i$.

(iv) and (v) follow from (ii) and (iii) by taking complements via De~Morgan's laws.
\end{proof}

\begin{proposition}[Sequential Characterization of Closed Sets]
\label{prop:closed_seq}
A set $F \subseteq \R$ is closed if and only if it contains all its limit points:
whenever $(x_n)$ is a sequence in $F$ with $x_n \to L$, we have $L \in F$.
\end{proposition}

\begin{proof}
$(\Rightarrow)$ Suppose $F$ is closed, $(x_n) \subseteq F$, and $x_n \to L$.
If $L \notin F$, then $L \in \R \setminus F$ which is open, so there exists
$\varepsilon > 0$ with $(L - \varepsilon, L + \varepsilon) \subseteq \R \setminus F$.
But $x_n \to L$ gives $x_n \in (L - \varepsilon, L + \varepsilon)$ for large $n$,
contradicting $x_n \in F$.

$(\Leftarrow)$ Suppose $F$ contains all its limit points. We show $\R \setminus F$
is open. Let $x \in \R \setminus F$. If no $\varepsilon$-neighborhood of $x$ is
contained in $\R \setminus F$, then for each $n$ there exists $x_n \in F$ with
$|x_n - x| < 1/n$, giving $x_n \to x$ with $(x_n) \subseteq F$, so $x \in F$,
a contradiction.
\end{proof}

\begin{definition}[Closure, Interior, Boundary]
\label{def:closure_interior}
Let $S \subseteq \R$.
\begin{enumerate}[nosep,label=(\roman*)]
    \item The \emph{closure} $\overline{S}$ is the smallest closed set containing
    $S$: $\overline{S} = \bigcap\{F : F \text{ closed}, S \subseteq F\}$.
    Equivalently, $x \in \overline{S}$ iff every neighborhood of $x$ intersects $S$.
    \item The \emph{interior} $S^\circ$ is the largest open set contained in $S$:
    $S^\circ = \bigcup\{U : U \text{ open}, U \subseteq S\}$.
    \item The \emph{boundary} $\partial S = \overline{S} \setminus S^\circ$.
\end{enumerate}
\end{definition}

\begin{example}
\label{ex:closure}
$\overline{\Q} = \R$ and $\Q^\circ = \emptyset$ (by density of both rationals and
irrationals). Thus $\partial \Q = \R$.
\end{example}

\subsection{Compactness and the Heine--Borel Theorem}
\label{subsec:compactness}

\begin{definition}[Open Cover, Compactness]
\label{def:compact}
An \emph{open cover} of $S \subseteq \R$ is a collection $\{U_\alpha\}$ of open
sets with $S \subseteq \bigcup_\alpha U_\alpha$. A set $K \subseteq \R$ is
\emph{compact} if every open cover of $K$ has a finite subcover.
\end{definition}

\begin{theorem}[Heine--Borel]
\label{thm:heine_borel}
A subset $K \subseteq \R$ is compact if and only if it is closed and bounded.
\end{theorem}

\begin{proof}
$(\Rightarrow)$ Suppose $K$ is compact.

\emph{$K$ is bounded:} The collection $\{(-n, n) : n \in \N\}$ is an open cover
of $K$. By compactness, finitely many of these cover $K$, say
$K \subseteq (-N, N)$. So $K$ is bounded.

\emph{$K$ is closed:} Let $x \notin K$. For each $y \in K$, let
$r_y = |x - y|/2 > 0$. The open sets $U_y = (y - r_y, y + r_y)$ cover $K$.
By compactness, $K \subseteq U_{y_1} \cup \cdots \cup U_{y_m}$ for some finite
collection. Let $\delta = \min(r_{y_1}, \ldots, r_{y_m}) > 0$. Then
$(x - \delta, x + \delta) \cap K = \emptyset$ (since any $z$ in this
neighborhood satisfies $|z - x| < \delta \le r_{y_i}$, hence
$|z - y_i| \ge |x - y_i| - |x - z| > 2r_{y_i} - r_{y_i} = r_{y_i}$,
so $z \notin U_{y_i}$ for each $i$). Thus $\R \setminus K$ is open, so $K$ is
closed.

$(\Leftarrow)$ Suppose $K$ is closed and bounded, say $K \subseteq [a, b]$. Let
$\{U_\alpha\}$ be an open cover of $K$. We show it has a finite subcover by
contradiction.

Bisect $[a, b]$ into $[a, (a+b)/2]$ and $[(a+b)/2, b]$. At least one half, call
it $[a_1, b_1]$, has the property that $K \cap [a_1, b_1]$ is not covered by any
finite subcollection of $\{U_\alpha\}$. Continue bisecting: at step $n$, we have
$[a_n, b_n]$ with $b_n - a_n = (b-a)/2^n$ such that $K \cap [a_n, b_n]$ cannot
be finitely covered.

By the nested intervals property (Theorem~\ref{thm:nested_intervals}), there
exists $x \in \bigcap_{n=1}^\infty [a_n, b_n]$. Since each $[a_n, b_n]$ intersects
$K$ nontrivially (otherwise a finite --- in fact, empty --- subcover would exist),
choose $x_n \in K \cap [a_n, b_n]$. Then $|x_n - x| \le b_n - a_n = (b-a)/2^n
\to 0$, so $x_n \to x$. Since $K$ is closed, $x \in K$.

Since $\{U_\alpha\}$ covers $K$, there exists $U_\beta$ with $x \in U_\beta$.
Since $U_\beta$ is open, there exists $\varepsilon > 0$ with
$(x - \varepsilon, x + \varepsilon) \subseteq U_\beta$. For $n$ large enough that
$(b-a)/2^n < \varepsilon$, we have $[a_n, b_n] \subseteq (x - \varepsilon,
x + \varepsilon) \subseteq U_\beta$. Then $\{U_\beta\}$ is a finite subcover
of $K \cap [a_n, b_n]$, a contradiction.
\end{proof}

\begin{proposition}[Properties of Compact Sets]
\label{prop:compact_props}
\begin{enumerate}[nosep,label=(\roman*)]
    \item A closed subset of a compact set is compact.
    \item A compact set is sequentially compact: every sequence in $K$ has a
    subsequence converging to a point in $K$.
\end{enumerate}
\end{proposition}

\begin{proof}
(i) Let $F \subseteq K$ be closed with $K$ compact. If $\{U_\alpha\}$ is an open
cover of $F$, then $\{U_\alpha\} \cup \{\R \setminus F\}$ is an open cover of $K$
(since $\R \setminus F$ is open). By compactness of $K$, a finite subcover exists.
Removing $\R \setminus F$ (if present) leaves a finite subcover of $F$.

(ii) If $(x_n)$ is a sequence in $K$, then $(x_n)$ is bounded (since $K$ is bounded
by Heine--Borel). By Bolzano--Weierstrass, a subsequence converges to some $L$.
Since $K$ is closed, $L \in K$.
\end{proof}

\begin{rlconnection}[Compactness Beyond $\R^n$]
Compactness is used throughout measure theory and functional analysis. The
Heine--Borel theorem characterizes compact subsets of $\R^n$. In functional
analysis (Phase~01), compactness in infinite-dimensional spaces is more subtle
and requires the Arzel\`{a}--Ascoli theorem.
\end{rlconnection}

%% ============================================================
%% SECTION 5: CONTINUITY
%% ============================================================
\section{Continuity}
\label{sec:continuity}

\subsection{Definition and Sequential Characterization}
\label{subsec:cont_def}

\begin{definition}[Continuity]
\label{def:continuity}
Let $f : A \to \R$ with $A \subseteq \R$ and let $c \in A$. We say $f$ is
\emph{continuous at $c$} if for every $\varepsilon > 0$ there exists $\delta > 0$
such that $|f(x) - f(c)| < \varepsilon$ whenever $x \in A$ and $|x - c| < \delta$.
We say $f$ is \emph{continuous on $A$} if it is continuous at every $c \in A$.
\end{definition}

\begin{theorem}[Sequential Characterization of Continuity]
\label{thm:seq_continuity}
Let $f : A \to \R$ and $c \in A$. Then $f$ is continuous at $c$ if and only if
for every sequence $(x_n)$ in $A$ with $x_n \to c$, we have $f(x_n) \to f(c)$.
\end{theorem}

\begin{proof}
$(\Rightarrow)$ Suppose $f$ is continuous at $c$ and $x_n \to c$. Given
$\varepsilon > 0$, choose $\delta > 0$ from the definition of continuity, then
choose $N$ with $|x_n - c| < \delta$ for $n \ge N$. Then
$|f(x_n) - f(c)| < \varepsilon$ for $n \ge N$.

$(\Leftarrow)$ We prove the contrapositive. Suppose $f$ is not continuous at $c$.
Then there exists $\varepsilon_0 > 0$ such that for every $\delta > 0$ there exists
$x \in A$ with $|x - c| < \delta$ but $|f(x) - f(c)| \ge \varepsilon_0$. Taking
$\delta = 1/n$, we obtain a sequence $x_n \in A$ with $|x_n - c| < 1/n$ (so
$x_n \to c$) but $|f(x_n) - f(c)| \ge \varepsilon_0$ (so $f(x_n) \not\to f(c)$).
\end{proof}

\begin{proposition}[Algebraic Properties of Continuous Functions]
\label{prop:cont_algebra}
If $f, g : A \to \R$ are continuous at $c \in A$, then so are $f + g$, $fg$,
and $f/g$ (the last provided $g(c) \neq 0$). If $g : A \to \R$ is continuous
at $c$ and $h : B \to \R$ is continuous at $g(c)$ with $g(A) \subseteq B$,
then $h \circ g$ is continuous at $c$.
\end{proposition}

\begin{proof}
The first three assertions follow immediately from Theorem~\ref{thm:seq_continuity}
and the algebraic limit theorems (Theorem~\ref{thm:alg_limit}). For composition:
if $x_n \to c$, then $g(x_n) \to g(c)$ by continuity of $g$, then
$h(g(x_n)) \to h(g(c))$ by continuity of $h$.
\end{proof}

\subsection{Extreme Value Theorem and Intermediate Value Theorem}
\label{subsec:evt_ivt}

\begin{theorem}[Extreme Value Theorem]
\label{thm:evt}
If $f : K \to \R$ is continuous and $K$ is compact, then $f$ attains its maximum
and minimum: there exist $x_*, x^* \in K$ with $f(x_*) \le f(x) \le f(x^*)$ for
all $x \in K$.
\end{theorem}

\begin{proof}
We show $f$ attains its maximum; the minimum argument is analogous (apply the
result to $-f$).

\emph{Step 1: $f$ is bounded above.} Suppose not. Then for each $n$ there exists
$x_n \in K$ with $f(x_n) > n$. By sequential compactness
(Proposition~\ref{prop:compact_props}(ii)), $(x_n)$ has a subsequence
$x_{n_k} \to c \in K$. By continuity, $f(x_{n_k}) \to f(c)$, but
$f(x_{n_k}) > n_k \to \infty$, a contradiction.

\emph{Step 2: The supremum is attained.} Let $M = \sup\{f(x) : x \in K\}$, which
is finite by Step~1. By the characterization of supremum, for each $n$ there
exists $x_n \in K$ with $f(x_n) > M - 1/n$. By sequential compactness,
$x_{n_k} \to x^* \in K$ for some subsequence. By continuity,
$f(x^*) = \lim f(x_{n_k}) = M$.
\end{proof}

\begin{theorem}[Intermediate Value Theorem]
\label{thm:ivt}
If $f : [a, b] \to \R$ is continuous and $f(a) < v < f(b)$ (or $f(a) > v > f(b)$),
then there exists $c \in (a, b)$ with $f(c) = v$.
\end{theorem}

\begin{proof}
Assume $f(a) < v < f(b)$. Let $S = \{x \in [a, b] : f(x) < v\}$. Then $S$ is
nonempty (since $a \in S$) and bounded above (by $b$), so
$c = \sup S$ exists with $a \le c \le b$.

\emph{Claim: $f(c) = v$.} Choose a sequence $(s_n)$ in $S$ with $s_n \to c$.
By continuity, $f(c) = \lim f(s_n) \le v$ (since $f(s_n) < v$). Now if $f(c) < v$,
then by continuity there exists $\delta > 0$ with $f(x) < v$ for
$x \in (c - \delta, c + \delta) \cap [a, b]$. Since $c < b$ (because $f(c) < v < f(b)$),
there exists $x \in (c, c + \delta) \cap [a, b]$ with $f(x) < v$, contradicting
$c = \sup S$. Therefore $f(c) = v$.
\end{proof}

\subsection{Uniform Continuity}
\label{subsec:uniform_cont}

\begin{definition}[Uniform Continuity]
\label{def:uniform_cont}
A function $f : A \to \R$ is \emph{uniformly continuous} on $A$ if for every
$\varepsilon > 0$ there exists $\delta > 0$ such that $|f(x) - f(y)| < \varepsilon$
whenever $x, y \in A$ with $|x - y| < \delta$.
\end{definition}

\begin{remark}
\label{rem:unif_vs_cont}
The distinction from ordinary continuity: in uniform continuity, $\delta$ depends
only on $\varepsilon$, not on the point. Ordinary continuity allows $\delta$ to vary
with the point $c$.
\end{remark}

\begin{theorem}[Continuous on Compact $\Rightarrow$ Uniformly Continuous]
\label{thm:compact_unif_cont}
If $K \subseteq \R$ is compact and $f : K \to \R$ is continuous, then $f$ is
uniformly continuous on $K$.
\end{theorem}

\begin{proof}
Suppose $f$ is not uniformly continuous on $K$. Then there exists
$\varepsilon_0 > 0$ such that for every $n \in \N$ there exist $x_n, y_n \in K$
with $|x_n - y_n| < 1/n$ but $|f(x_n) - f(y_n)| \ge \varepsilon_0$.

By sequential compactness, $(x_n)$ has a subsequence $x_{n_k} \to c \in K$.
Since $|y_{n_k} - c| \le |y_{n_k} - x_{n_k}| + |x_{n_k} - c| < 1/n_k +
|x_{n_k} - c| \to 0$, we also have $y_{n_k} \to c$. By continuity of $f$ at $c$:
\[
|f(x_{n_k}) - f(y_{n_k})| \le |f(x_{n_k}) - f(c)| + |f(c) - f(y_{n_k})| \to 0,
\]
contradicting $|f(x_{n_k}) - f(y_{n_k})| \ge \varepsilon_0$.
\end{proof}

%% ============================================================
%% SECTION 6: DIFFERENTIATION
%% ============================================================
\section{Differentiation}
\label{sec:differentiation}

\subsection{Definition and Basic Properties}
\label{subsec:diff_def}

\begin{definition}[Derivative]
\label{def:derivative}
Let $f : (a, b) \to \R$ and $c \in (a, b)$. The \emph{derivative of $f$ at $c$}
is
\[
f'(c) = \lim_{h \to 0} \frac{f(c + h) - f(c)}{h},
\]
provided this limit exists. If $f'(c)$ exists, we say $f$ is \emph{differentiable
at $c$}.
\end{definition}

\begin{theorem}[Differentiable Implies Continuous]
\label{thm:diff_cont}
If $f$ is differentiable at $c$, then $f$ is continuous at $c$.
\end{theorem}

\begin{proof}
We have
\[
f(c + h) - f(c) = \frac{f(c+h) - f(c)}{h} \cdot h \to f'(c) \cdot 0 = 0
\]
as $h \to 0$, so $f(c+h) \to f(c)$.
\end{proof}

\begin{proposition}[Algebraic Differentiation Rules]
\label{prop:diff_rules}
If $f$ and $g$ are differentiable at $c$, then:
\begin{enumerate}[nosep,label=(\roman*)]
    \item $(f + g)'(c) = f'(c) + g'(c)$.
    \item (Product rule) $(fg)'(c) = f'(c)g(c) + f(c)g'(c)$.
    \item (Quotient rule) If $g(c) \neq 0$, then
    $(f/g)'(c) = \bigl(f'(c)g(c) - f(c)g'(c)\bigr)/g(c)^2$.
\end{enumerate}
\end{proposition}

\begin{proof}
We prove~(ii). Write
\begin{align*}
f(c+h)g(c+h) - f(c)g(c)
&= f(c+h)g(c+h) - f(c)g(c+h) + f(c)g(c+h) - f(c)g(c) \\
&= \bigl(f(c+h) - f(c)\bigr)g(c+h) + f(c)\bigl(g(c+h) - g(c)\bigr).
\end{align*}
Dividing by $h$ and taking $h \to 0$, using $g(c+h) \to g(c)$ (since $g$ is
continuous at $c$ by Theorem~\ref{thm:diff_cont}), gives
$f'(c)g(c) + f(c)g'(c)$.
\end{proof}

\begin{theorem}[Chain Rule]
\label{thm:chain_rule}
Let $g$ be differentiable at $c$ and $f$ be differentiable at $g(c)$. Then
$f \circ g$ is differentiable at $c$ and $(f \circ g)'(c) = f'(g(c)) \cdot g'(c)$.
\end{theorem}

\begin{proof}
Define the auxiliary function
\[
\varphi(y) =
\begin{cases}
\dfrac{f(y) - f(g(c))}{y - g(c)} & \text{if } y \neq g(c), \\[6pt]
f'(g(c)) & \text{if } y = g(c).
\end{cases}
\]
Since $f$ is differentiable at $g(c)$, $\varphi$ is continuous at $g(c)$.
For $h \neq 0$ with $g(c+h) \neq g(c)$:
\[
\frac{f(g(c+h)) - f(g(c))}{h}
= \varphi(g(c+h)) \cdot \frac{g(c+h) - g(c)}{h}.
\]
This identity also holds when $g(c+h) = g(c)$ (both sides are zero, using the
convention that $\varphi(g(c)) = f'(g(c))$ and $g(c+h) - g(c) = 0$).
Taking $h \to 0$: $\varphi(g(c+h)) \to \varphi(g(c)) = f'(g(c))$ (by continuity
of $\varphi$ at $g(c)$ and continuity of $g$ at $c$), and
$(g(c+h) - g(c))/h \to g'(c)$. Therefore
$(f \circ g)'(c) = f'(g(c)) \cdot g'(c)$.
\end{proof}

\subsection{Mean Value Theorem and Consequences}
\label{subsec:mvt}

\begin{theorem}[Fermat's Theorem on Local Extrema]
\label{thm:fermat}
If $f : (a, b) \to \R$ is differentiable at $c \in (a, b)$ and $f$ has a local
extremum at $c$, then $f'(c) = 0$.
\end{theorem}

\begin{proof}
Suppose $f$ has a local maximum at $c$ (the minimum case is analogous). There
exists $\delta > 0$ with $f(x) \le f(c)$ for $|x - c| < \delta$. For
$0 < h < \delta$:
\[
\frac{f(c+h) - f(c)}{h} \le 0,
\]
so $f'(c) = \lim_{h \to 0^+} \frac{f(c+h)-f(c)}{h} \le 0$. For
$-\delta < h < 0$:
\[
\frac{f(c+h) - f(c)}{h} \ge 0,
\]
so $f'(c) = \lim_{h \to 0^-} \frac{f(c+h)-f(c)}{h} \ge 0$. Together, $f'(c) = 0$.
\end{proof}

\begin{theorem}[Rolle's Theorem]
\label{thm:rolle}
If $f : [a, b] \to \R$ is continuous on $[a, b]$, differentiable on $(a, b)$,
and $f(a) = f(b)$, then there exists $c \in (a, b)$ with $f'(c) = 0$.
\end{theorem}

\begin{proof}
If $f$ is constant on $[a, b]$, then $f'(c) = 0$ for all $c \in (a, b)$.
Otherwise, $f$ is not identically $f(a)$, so $f$ attains a maximum or minimum
value different from $f(a) = f(b)$. By the extreme value theorem
(Theorem~\ref{thm:evt}), this extremum is attained at some $c \in [a, b]$.
Since the extreme value differs from $f(a) = f(b)$, the point $c$ cannot be
an endpoint, so $c \in (a, b)$. By Fermat's theorem (Theorem~\ref{thm:fermat}),
$f'(c) = 0$.
\end{proof}

\begin{theorem}[Mean Value Theorem]
\label{thm:mvt}
If $f : [a, b] \to \R$ is continuous on $[a, b]$ and differentiable on $(a, b)$,
then there exists $c \in (a, b)$ with
\[
f'(c) = \frac{f(b) - f(a)}{b - a}.
\]
\end{theorem}

\begin{proof}
Define $g(x) = f(x) - \frac{f(b) - f(a)}{b - a}(x - a)$. Then $g$ is continuous
on $[a, b]$, differentiable on $(a, b)$, and $g(a) = f(a) = g(b)$ (verify:
$g(b) = f(b) - (f(b) - f(a)) = f(a)$). By Rolle's theorem, there exists
$c \in (a, b)$ with $0 = g'(c) = f'(c) - \frac{f(b) - f(a)}{b - a}$.
\end{proof}

\subsection{Taylor's Theorem}
\label{subsec:taylor}

\begin{theorem}[Taylor's Theorem with Lagrange Remainder]
\label{thm:taylor}
Let $f : [a, b] \to \R$ have continuous derivatives $f', f'', \ldots, f^{(n)}$
on $[a, b]$ and let $f^{(n+1)}$ exist on $(a, b)$. Then for any $x \in [a, b]$
and $c \in [a, b]$:
\[
f(x) = \sum_{k=0}^{n} \frac{f^{(k)}(c)}{k!}(x - c)^k
+ \frac{f^{(n+1)}(\xi)}{(n+1)!}(x - c)^{n+1}
\]
for some $\xi$ between $c$ and $x$.
\end{theorem}

\begin{proof}
We may assume $x \neq c$; the result is trivial otherwise. Define the remainder
$R = f(x) - \sum_{k=0}^{n} \frac{f^{(k)}(c)}{k!}(x-c)^k$ and choose $M$ so that
\[
R = \frac{M}{(n+1)!}(x - c)^{n+1}.
\]
We need to show $M = f^{(n+1)}(\xi)$ for some $\xi$ between $c$ and $x$. Define
\[
g(t) = f(x) - \sum_{k=0}^{n} \frac{f^{(k)}(t)}{k!}(x - t)^k
- \frac{M}{(n+1)!}(x - t)^{n+1}.
\]
Then $g(c) = f(x) - \sum_{k=0}^n \frac{f^{(k)}(c)}{k!}(x-c)^k -
\frac{M}{(n+1)!}(x-c)^{n+1} = R - R = 0$, and $g(x) = f(x) - f(x) = 0$.
By Rolle's theorem, there exists $\xi$ between $c$ and $x$ with $g'(\xi) = 0$.

Computing $g'(t)$ by differentiating each term and using telescoping cancellation:
\begin{align*}
g'(t) &= -\sum_{k=0}^{n} \left[\frac{f^{(k+1)}(t)}{k!}(x-t)^k
- \frac{f^{(k)}(t)}{(k-1)!}(x-t)^{k-1}\right]
+ \frac{M}{n!}(x-t)^n,
\end{align*}
where the sum telescopes (the second part of the $k$-th term cancels the first
part of the $(k-1)$-th term), leaving
\[
g'(t) = -\frac{f^{(n+1)}(t)}{n!}(x - t)^n + \frac{M}{n!}(x-t)^n
= \frac{(x-t)^n}{n!}\bigl(M - f^{(n+1)}(t)\bigr).
\]
Setting $g'(\xi) = 0$ and noting $(x - \xi)^n \neq 0$ (since $\xi$ is between
$c$ and $x$, and $c \neq x$) gives $M = f^{(n+1)}(\xi)$.
\end{proof}

\begin{rlconnection}[Taylor's Theorem in Optimization and Measure Theory]
Taylor's theorem is essential for optimization theory (Phase~01). Gradient descent
convergence proofs use the first-order Taylor expansion with remainder to bound
function value changes. The mean value theorem appears in many
$\varepsilon$-$\delta$ arguments throughout measure theory.
\end{rlconnection}

%% ============================================================
%% SECTION 7: RIEMANN INTEGRATION
%% ============================================================
\section{Riemann Integration}
\label{sec:riemann}

\subsection{Partitions and Riemann Sums}
\label{subsec:partitions}

\begin{definition}[Partition]
\label{def:partition}
A \emph{partition} of $[a, b]$ is a finite set $P = \{x_0, x_1, \ldots, x_n\}$
with $a = x_0 < x_1 < \cdots < x_n = b$.
\end{definition}

\begin{definition}[Upper and Lower Sums]
\label{def:upper_lower_sums}
Let $f : [a, b] \to \R$ be bounded and $P = \{x_0, \ldots, x_n\}$ a partition.
Define
\[
M_k = \sup_{x \in [x_{k-1}, x_k]} f(x), \qquad
m_k = \inf_{x \in [x_{k-1}, x_k]} f(x), \qquad
\Delta x_k = x_k - x_{k-1}.
\]
The \emph{upper sum} and \emph{lower sum} are
\[
U(f, P) = \sum_{k=1}^n M_k \,\Delta x_k, \qquad
L(f, P) = \sum_{k=1}^n m_k \,\Delta x_k.
\]
\end{definition}

\begin{lemma}[Refinement Lemma]
\label{lem:refinement}
If $Q$ is a refinement of $P$ (i.e., $P \subseteq Q$), then
$L(f, P) \le L(f, Q) \le U(f, Q) \le U(f, P)$.
\end{lemma}

\begin{proof}
It suffices to show this when $Q$ has one more point than $P$. Say
$Q = P \cup \{x^*\}$ with $x_{k-1} < x^* < x_k$. On $[x_{k-1}, x_k]$,
the lower sum contribution changes from $m_k(x_k - x_{k-1})$ to
$m'(x^* - x_{k-1}) + m''(x_k - x^*)$ where
$m' = \inf_{[x_{k-1}, x^*]} f \ge m_k$ and
$m'' = \inf_{[x^*, x_k]} f \ge m_k$. Thus the lower sum does not decrease.
The upper sum argument is analogous (infima become suprema, inequality reverses).
The middle inequality $L(f, Q) \le U(f, Q)$ is immediate from $m_k \le M_k$.
\end{proof}

\begin{corollary}
\label{cor:lower_le_upper}
For any two partitions $P, Q$: $L(f, P) \le U(f, Q)$.
\end{corollary}

\begin{proof}
$L(f, P) \le L(f, P \cup Q) \le U(f, P \cup Q) \le U(f, Q)$.
\end{proof}

\begin{definition}[Riemann Integrability]
\label{def:riemann_int}
A bounded function $f : [a, b] \to \R$ is \emph{Riemann integrable} if
\[
\sup_P L(f, P) = \inf_P U(f, P),
\]
where the supremum and infimum are over all partitions $P$ of $[a, b]$.
This common value is the \emph{Riemann integral} $\int_a^b f(x)\,dx$.
By Corollary~\ref{cor:lower_le_upper}, the sup and inf exist and satisfy
$\sup_P L(f,P) \le \inf_P U(f,P)$.
\end{definition}

\begin{remark}[Riemann Criterion]
\label{rem:riemann_criterion}
Equivalently, $f$ is Riemann integrable if and only if for every $\varepsilon > 0$
there exists a partition $P$ with $U(f, P) - L(f, P) < \varepsilon$.
\end{remark}

\subsection{Continuous Functions Are Riemann Integrable}
\label{subsec:cont_integrable}

\begin{theorem}
\label{thm:cont_integrable}
If $f : [a, b] \to \R$ is continuous, then $f$ is Riemann integrable.
\end{theorem}

\begin{proof}
Since $[a, b]$ is compact, $f$ is uniformly continuous
(Theorem~\ref{thm:compact_unif_cont}). Given $\varepsilon > 0$, choose $\delta > 0$
with $|f(x) - f(y)| < \varepsilon/(b - a)$ whenever $|x - y| < \delta$. Let $P$
be any partition with mesh $\|P\| = \max_k \Delta x_k < \delta$. On each subinterval
$[x_{k-1}, x_k]$, continuity on the compact interval gives
$M_k - m_k = f(u_k) - f(v_k)$ for some $u_k, v_k \in [x_{k-1}, x_k]$
(extreme value theorem), and $|u_k - v_k| \le \Delta x_k < \delta$, so
$M_k - m_k < \varepsilon/(b-a)$. Therefore:
\[
U(f, P) - L(f, P) = \sum_{k=1}^n (M_k - m_k)\Delta x_k
< \frac{\varepsilon}{b - a} \sum_{k=1}^n \Delta x_k = \varepsilon.
\]
By the Riemann criterion (Remark~\ref{rem:riemann_criterion}), $f$ is integrable.
\end{proof}

\subsection{Properties of the Riemann Integral}
\label{subsec:riemann_props}

\begin{proposition}[Properties]
\label{prop:riemann_props}
Let $f, g : [a, b] \to \R$ be Riemann integrable. Then:
\begin{enumerate}[nosep,label=(\roman*)]
    \item \textbf{Linearity:} $\int_a^b (\alpha f + \beta g) = \alpha \int_a^b f
    + \beta \int_a^b g$ for all $\alpha, \beta \in \R$.
    \item \textbf{Monotonicity:} If $f(x) \le g(x)$ for all $x \in [a, b]$, then
    $\int_a^b f \le \int_a^b g$.
    \item \textbf{Additivity:} For $a < c < b$, $f$ is integrable on $[a, c]$ and
    $[c, b]$, and $\int_a^b f = \int_a^c f + \int_c^b f$.
    \item \textbf{Triangle inequality:} $|f|$ is integrable and
    $|\int_a^b f| \le \int_a^b |f|$.
\end{enumerate}
\end{proposition}

\begin{proof}
We prove~(ii). For any partition $P$,
$\inf_{[x_{k-1}, x_k]} f \le \inf_{[x_{k-1}, x_k]} g$ is \emph{not} generally
true; rather we use: $\sup_P L(f, P) \le \sup_P L(g, P)$ since for each partition,
$m_k^f \le m_k^g$ implies $L(f, P) \le L(g, P)$. Wait --- this is incorrect since
$m_k^f \le m_k^g$ does \emph{not} follow from $f \le g$. Actually it does: if
$f(x) \le g(x)$ for all $x$, then $\inf_{[x_{k-1},x_k]} f(x) \le
\inf_{[x_{k-1},x_k]} g(x)$. Therefore $L(f,P) \le L(g,P)$ for all $P$, and
$\int_a^b f = \sup_P L(f,P) \le \sup_P L(g,P) = \int_a^b g$.

We prove (i) for the case $\alpha = 1, \beta = 1$ (the general case follows
similarly). Given $\varepsilon > 0$, choose partitions $P_1, P_2$ with
$U(f, P_1) - L(f, P_1) < \varepsilon/2$ and $U(g, P_2) - L(g, P_2) < \varepsilon/2$.
Let $P = P_1 \cup P_2$. On each subinterval of $P$,
$\sup(f + g) \le \sup f + \sup g$ and $\inf(f + g) \ge \inf f + \inf g$. Therefore
$U(f+g, P) \le U(f, P) + U(g, P)$ and $L(f+g, P) \ge L(f, P) + L(g, P)$.
Thus:
\[
U(f+g, P) - L(f+g, P) \le \bigl(U(f,P) - L(f,P)\bigr) + \bigl(U(g,P) - L(g,P)\bigr)
< \varepsilon,
\]
so $f + g$ is integrable. Moreover, $L(f,P) + L(g,P) \le L(f+g,P) \le
\int_a^b(f+g) \le U(f+g,P) \le U(f,P) + U(g,P)$. Since
$\int_a^b f + \int_a^b g$ also lies between $L(f,P) + L(g,P)$ and
$U(f,P) + U(g,P)$, the two values differ by at most $\varepsilon$.
As $\varepsilon$ is arbitrary, equality holds.
\end{proof}

\subsection{Fundamental Theorem of Calculus}
\label{subsec:ftc}

\begin{theorem}[Fundamental Theorem of Calculus, Part I]
\label{thm:ftc1}
Let $f : [a, b] \to \R$ be Riemann integrable. Define
$F(x) = \int_a^x f(t)\,dt$ for $x \in [a, b]$. Then $F$ is continuous on $[a, b]$.
If $f$ is continuous at $c \in (a, b)$, then $F$ is differentiable at $c$ and
$F'(c) = f(c)$.
\end{theorem}

\begin{proof}
\emph{Continuity:} Since $f$ is bounded, say $|f(t)| \le M$, for $x, y \in [a, b]$
with $x < y$:
\[
|F(y) - F(x)| = \Bigl|\int_x^y f(t)\,dt\Bigr| \le M(y - x).
\]
So $|F(y) - F(x)| \le M|y - x|$, and $F$ is (Lipschitz) continuous.

\emph{Differentiability:} Suppose $f$ is continuous at $c$. Given
$\varepsilon > 0$, choose $\delta > 0$ with $|f(t) - f(c)| < \varepsilon$ for
$|t - c| < \delta$. For $0 < |h| < \delta$:
\[
\frac{F(c+h) - F(c)}{h} - f(c)
= \frac{1}{h}\int_c^{c+h} f(t)\,dt - f(c)
= \frac{1}{h}\int_c^{c+h} \bigl(f(t) - f(c)\bigr)\,dt.
\]
Therefore
\[
\Bigl|\frac{F(c+h) - F(c)}{h} - f(c)\Bigr|
\le \frac{1}{|h|}\int_{\min(c,c+h)}^{\max(c,c+h)} |f(t) - f(c)|\,dt
< \frac{1}{|h|} \cdot \varepsilon \cdot |h| = \varepsilon. \qedhere
\]
\end{proof}

\begin{theorem}[Fundamental Theorem of Calculus, Part II]
\label{thm:ftc2}
If $f : [a, b] \to \R$ is Riemann integrable and $F : [a, b] \to \R$ is continuous
on $[a, b]$, differentiable on $(a, b)$, with $F'(x) = f(x)$ for all
$x \in (a, b)$, then
\[
\int_a^b f(x)\,dx = F(b) - F(a).
\]
\end{theorem}

\begin{proof}
Let $P = \{x_0, x_1, \ldots, x_n\}$ be any partition of $[a, b]$. By the mean
value theorem (Theorem~\ref{thm:mvt}), on each subinterval there exists
$c_k \in (x_{k-1}, x_k)$ with $F(x_k) - F(x_{k-1}) = F'(c_k)\Delta x_k =
f(c_k)\Delta x_k$. Therefore:
\[
F(b) - F(a) = \sum_{k=1}^n \bigl(F(x_k) - F(x_{k-1})\bigr)
= \sum_{k=1}^n f(c_k)\Delta x_k.
\]
Since $m_k \le f(c_k) \le M_k$, we get
$L(f, P) \le F(b) - F(a) \le U(f, P)$ for every partition $P$. Since $f$ is
integrable, the common value $\int_a^b f$ is the unique number between all lower
and upper sums, so $\int_a^b f = F(b) - F(a)$.
\end{proof}

\begin{warning}[Limitations of the Riemann Integral]
The Riemann integral cannot handle many natural limits. For example, let
$q_1, q_2, q_3, \ldots$ be an enumeration of the rationals in $[0,1]$ and define
$f_n = \mathbf{1}_{\{q_1, \ldots, q_n\}}$ (the indicator of the first $n$
rationals). Each $f_n$ is Riemann integrable with $\int_0^1 f_n = 0$, but the
pointwise limit is the Dirichlet function $\mathbf{1}_\Q$, which is \textbf{not}
Riemann integrable. This failure motivates the Lebesgue integral in Phase~01.
\end{warning}

%% ============================================================
%% SECTION 8: SEQUENCES AND SERIES OF FUNCTIONS
%% ============================================================
\section{Sequences and Series of Functions}
\label{sec:function_sequences}

\subsection{Pointwise and Uniform Convergence}
\label{subsec:pointwise_uniform}

\begin{definition}[Pointwise Convergence]
\label{def:pointwise}
A sequence of functions $f_n : A \to \R$ \emph{converges pointwise} to
$f : A \to \R$ if for every $x \in A$, $\lim_{n \to \infty} f_n(x) = f(x)$.
That is, for every $x \in A$ and every $\varepsilon > 0$, there exists
$N = N(x, \varepsilon) \in \N$ such that $|f_n(x) - f(x)| < \varepsilon$
for all $n \ge N$.
\end{definition}

\begin{definition}[Uniform Convergence]
\label{def:uniform}
A sequence $f_n : A \to \R$ \emph{converges uniformly} to $f : A \to \R$ if
for every $\varepsilon > 0$ there exists $N = N(\varepsilon) \in \N$ such that
$|f_n(x) - f(x)| < \varepsilon$ for all $n \ge N$ and all $x \in A$.
\end{definition}

\begin{remark}
\label{rem:uniform_sup}
Uniform convergence is equivalent to
$\sup_{x \in A} |f_n(x) - f(x)| \to 0$
as $n \to \infty$. The key difference from pointwise convergence is that $N$
does not depend on $x$.
\end{remark}

\begin{example}[Pointwise but Not Uniform Convergence]
\label{ex:pointwise_not_uniform}
Let $f_n : [0, 1] \to \R$ be defined by $f_n(x) = x^n$. For $x \in [0, 1)$,
$f_n(x) = x^n \to 0$; and $f_n(1) = 1$ for all $n$. So $f_n \to f$ pointwise,
where $f(x) = 0$ for $x \in [0, 1)$ and $f(1) = 1$. Each $f_n$ is continuous, but
the limit $f$ is discontinuous at $x = 1$. The convergence is not uniform:
$\sup_{x \in [0,1]} |f_n(x) - f(x)| \ge f_n(1 - 1/n) = (1 - 1/n)^n \to 1/e
\neq 0$.
\end{example}

\subsection{Interchange Theorems}
\label{subsec:interchange}

\begin{theorem}[Uniform Limit of Continuous Functions Is Continuous]
\label{thm:uniform_limit_cont}
Let $f_n : A \to \R$ be continuous for each $n$, and suppose $f_n \to f$
uniformly on $A$. Then $f$ is continuous on $A$.
\end{theorem}

\begin{proof}
Fix $c \in A$ and $\varepsilon > 0$. By uniform convergence, choose $N$ with
$|f_N(x) - f(x)| < \varepsilon/3$ for all $x \in A$. By continuity of $f_N$
at $c$, choose $\delta > 0$ with $|f_N(x) - f_N(c)| < \varepsilon/3$ for
$|x - c| < \delta$. Then for $|x - c| < \delta$:
\begin{align*}
|f(x) - f(c)|
&\le |f(x) - f_N(x)| + |f_N(x) - f_N(c)| + |f_N(c) - f(c)| \\
&< \frac{\varepsilon}{3} + \frac{\varepsilon}{3} + \frac{\varepsilon}{3}
= \varepsilon. \qedhere
\end{align*}
\end{proof}

\begin{theorem}[Interchange of Limit and Integral Under Uniform Convergence]
\label{thm:uniform_integral}
Let $f_n : [a, b] \to \R$ be Riemann integrable for each $n$, and suppose
$f_n \to f$ uniformly on $[a, b]$. Then $f$ is Riemann integrable and
\[
\int_a^b f(x)\,dx = \lim_{n \to \infty} \int_a^b f_n(x)\,dx.
\]
\end{theorem}

\begin{proof}
\emph{$f$ is integrable:} Given $\varepsilon > 0$, choose $N$ with
$|f_N(x) - f(x)| < \varepsilon/(4(b-a))$ for all $x$. Since $f_N$ is integrable,
choose a partition $P$ with $U(f_N, P) - L(f_N, P) < \varepsilon/2$. On each
subinterval, $\sup f \le \sup f_N + \varepsilon/(4(b-a))$ and
$\inf f \ge \inf f_N - \varepsilon/(4(b-a))$, so
\begin{align*}
U(f, P) - L(f, P)
&\le U(f_N, P) - L(f_N, P) + \frac{2\varepsilon}{4(b-a)} \cdot (b - a) \\
&< \frac{\varepsilon}{2} + \frac{\varepsilon}{2} = \varepsilon.
\end{align*}

\emph{Interchange:} For $n \ge N$ with $|f_n(x) - f(x)| < \varepsilon/(b-a)$:
\[
\Bigl|\int_a^b f_n - \int_a^b f\Bigr|
= \Bigl|\int_a^b (f_n - f)\Bigr|
\le \int_a^b |f_n - f|
< \frac{\varepsilon}{b - a} \cdot (b - a) = \varepsilon. \qedhere
\]
\end{proof}

\begin{theorem}[Weierstrass $M$-Test]
\label{thm:weierstrass_m}
Let $f_n : A \to \R$ be functions and suppose there exist constants $M_n \ge 0$
with $|f_n(x)| \le M_n$ for all $x \in A$ and $\sum_{n=1}^\infty M_n < \infty$.
Then $\sum_{n=1}^\infty f_n$ converges uniformly (and absolutely) on $A$.
\end{theorem}

\begin{proof}
For each $x$, $\sum |f_n(x)| \le \sum M_n < \infty$, so $\sum f_n(x)$ converges
absolutely. Let $S(x) = \sum_{n=1}^\infty f_n(x)$ and
$S_N(x) = \sum_{n=1}^N f_n(x)$. Then for all $x$:
\[
|S(x) - S_N(x)| = \Bigl|\sum_{n=N+1}^\infty f_n(x)\Bigr|
\le \sum_{n=N+1}^\infty |f_n(x)|
\le \sum_{n=N+1}^\infty M_n.
\]
The tail $\sum_{n=N+1}^\infty M_n \to 0$ as $N \to \infty$ (since $\sum M_n$
converges), and the bound is independent of $x$, so the convergence is uniform.
\end{proof}

\begin{keyresult}[From Uniform to Dominated Convergence]
The interchange theorems for uniform convergence are the starting point.
Measure theory (Phase~01) develops the far more powerful Monotone Convergence
Theorem and Dominated Convergence Theorem, which allow interchange of limit
and integral under much weaker conditions --- pointwise convergence suffices
when combined with a dominating function.
\end{keyresult}

%% ============================================================
%% SECTION 9: METRIC SPACES
%% ============================================================
\section{Metric Spaces}
\label{sec:metric_spaces}

\subsection{Definitions and Examples}
\label{subsec:metric_def}

\begin{definition}[Metric Space]
\label{def:metric_space}
A \emph{metric space} is a pair $(X, d)$ where $X$ is a nonempty set and
$d : X \times X \to [0, \infty)$ satisfies:
\begin{enumerate}[nosep,label=(\roman*)]
    \item $d(x, y) = 0$ if and only if $x = y$ \hfill (identity of indiscernibles),
    \item $d(x, y) = d(y, x)$ for all $x, y \in X$ \hfill (symmetry),
    \item $d(x, z) \le d(x, y) + d(y, z)$ for all $x, y, z \in X$
    \hfill (triangle inequality).
\end{enumerate}
\end{definition}

\begin{example}[Standard Examples]
\label{ex:metric_spaces}
\begin{enumerate}[nosep,label=(\alph*)]
    \item $(\R^n, d_2)$ where $d_2(x, y) = \bigl(\sum_{i=1}^n (x_i - y_i)^2
    \bigr)^{1/2}$ is the Euclidean metric.
    \item $(X, d_{\mathrm{disc}})$ where $d_{\mathrm{disc}}(x, y) = 0$ if $x = y$
    and $1$ if $x \neq y$ is the \emph{discrete metric} on any set $X$.
    \item $\bigl(C[a,b],\, d_\infty\bigr)$ where $C[a,b]$ is the set of
    continuous functions $f : [a, b] \to \R$ and $d_\infty(f, g) =
    \sup_{x \in [a,b]} |f(x) - g(x)|$ is the \emph{sup metric}
    (or uniform metric). This is well-defined because continuous functions on
    a compact set are bounded (Theorem~\ref{thm:evt}).
\end{enumerate}
\end{example}

\subsection{Topology, Convergence, and Completeness}
\label{subsec:metric_topology}

\begin{definition}[Open Ball, Open Set, Closed Set]
\label{def:metric_open}
In a metric space $(X, d)$:
\begin{enumerate}[nosep,label=(\roman*)]
    \item The \emph{open ball} of radius $r > 0$ centered at $x$ is
    $B(x, r) = \{y \in X : d(x, y) < r\}$.
    \item A set $U \subseteq X$ is \emph{open} if for every $x \in U$ there exists
    $r > 0$ with $B(x, r) \subseteq U$.
    \item A set $F \subseteq X$ is \emph{closed} if $X \setminus F$ is open.
\end{enumerate}
\end{definition}

\begin{definition}[Convergence in a Metric Space]
\label{def:metric_convergence}
A sequence $(x_n)$ in $(X, d)$ \emph{converges} to $x \in X$ if
$d(x_n, x) \to 0$, i.e., for every $\varepsilon > 0$ there exists $N$ with
$d(x_n, x) < \varepsilon$ for all $n \ge N$.
\end{definition}

\begin{definition}[Cauchy Sequence and Completeness]
\label{def:metric_complete}
A sequence $(x_n)$ in $(X, d)$ is \emph{Cauchy} if for every $\varepsilon > 0$
there exists $N$ with $d(x_n, x_m) < \varepsilon$ for all $n, m \ge N$.
A metric space is \emph{complete} if every Cauchy sequence converges (to an element
of $X$).
\end{definition}

\begin{example}
\label{ex:complete_spaces}
$\R$ is complete (Theorem~\ref{thm:cauchy_complete}). $\Q$ with the usual metric
is not complete: the sequence of decimal approximations to $\sqrt{2}$ is Cauchy
in $\Q$ but does not converge in $\Q$. The space $\bigl(C[a,b], d_\infty\bigr)$
is complete: this is a consequence of Theorem~\ref{thm:uniform_limit_cont}, since a
Cauchy sequence in the sup metric converges uniformly, and the uniform limit of
continuous functions is continuous.
\end{example}

\subsection{Compactness in Metric Spaces}
\label{subsec:metric_compact}

\begin{definition}[Sequential Compactness]
\label{def:seq_compact}
A metric space $(X, d)$ (or a subset $K \subseteq X$) is \emph{sequentially
compact} if every sequence in $K$ has a subsequence converging to a point in $K$.
\end{definition}

\begin{remark}
\label{rem:compact_equiv}
In a metric space, sequential compactness is equivalent to compactness (every open
cover has a finite subcover). The Heine--Borel theorem is a special case for $\R^n$:
a subset of $\R^n$ is (sequentially) compact if and only if it is closed and bounded.
\end{remark}

\subsection{The Contraction Mapping Theorem}
\label{subsec:contraction}

\begin{definition}[Contraction]
\label{def:contraction}
Let $(X, d)$ be a metric space. A map $T : X \to X$ is a \emph{contraction}
(or \emph{contraction mapping}) if there exists a constant $\gamma \in [0, 1)$
such that
\[
d(T(x), T(y)) \le \gamma \, d(x, y) \qquad \text{for all } x, y \in X.
\]
The constant $\gamma$ is called the \emph{contraction factor}.
\end{definition}

\begin{theorem}[Banach Fixed Point Theorem / Contraction Mapping Theorem]
\label{thm:banach_fixed_point}
Let $(X, d)$ be a complete metric space and $T : X \to X$ a contraction with
factor $\gamma \in [0, 1)$. Then:
\begin{enumerate}[nosep,label=(\roman*)]
    \item $T$ has a unique fixed point $x^* \in X$ (i.e., $T(x^*) = x^*$).
    \item For any $x_0 \in X$, the iterates $x_{n+1} = T(x_n)$ satisfy
    $x_n \to x^*$.
    \item The rate of convergence satisfies
    $d(x_n, x^*) \le \frac{\gamma^n}{1 - \gamma} d(x_0, x_1)$.
\end{enumerate}
\end{theorem}

\begin{proof}
\emph{Existence.}
Fix any $x_0 \in X$ and define $x_{n+1} = T(x_n)$. We show $(x_n)$ is Cauchy.
By the contraction property, $d(x_{n+1}, x_n) \le \gamma \, d(x_n, x_{n-1})$,
so by induction $d(x_{n+1}, x_n) \le \gamma^n d(x_1, x_0)$. For $m > n$:
\begin{align*}
d(x_m, x_n)
&\le \sum_{k=n}^{m-1} d(x_{k+1}, x_k)
\le \sum_{k=n}^{m-1} \gamma^k d(x_1, x_0)
\le \frac{\gamma^n}{1 - \gamma} d(x_1, x_0).
\end{align*}
Since $\gamma^n \to 0$, for any $\varepsilon > 0$ we can choose $N$ so that
$\gamma^N d(x_1, x_0)/(1 - \gamma) < \varepsilon$. Then $d(x_m, x_n) < \varepsilon$
for $m > n \ge N$, so $(x_n)$ is Cauchy. Since $X$ is complete, $x_n \to x^*$
for some $x^* \in X$.

We verify $T(x^*) = x^*$. By the triangle inequality and contraction:
\[
d(x^*, T(x^*)) \le d(x^*, x_{n+1}) + d(T(x_n), T(x^*))
\le d(x^*, x_{n+1}) + \gamma \, d(x_n, x^*) \to 0.
\]
So $d(x^*, T(x^*)) = 0$, giving $T(x^*) = x^*$.

\emph{Uniqueness.}
If $T(y^*) = y^*$ as well, then
$d(x^*, y^*) = d(T(x^*), T(y^*)) \le \gamma \, d(x^*, y^*)$.
Since $\gamma < 1$, this forces $d(x^*, y^*) = 0$, i.e., $x^* = y^*$.

\emph{Rate of convergence.}
Taking $m \to \infty$ in the estimate
$d(x_m, x_n) \le \gamma^n d(x_1, x_0)/(1 - \gamma)$ gives
$d(x^*, x_n) \le \gamma^n d(x_1, x_0)/(1 - \gamma)$.
\end{proof}

\begin{keyresult}[The Contraction Mapping Theorem in Reinforcement Learning]
This theorem is fundamental to reinforcement learning. The Bellman optimality
operator $T$ defined by
\[
(TV)(s) = \max_a \Bigl[R(s,a) + \gamma \sum_{s'} P(s' \mid s, a)\, V(s')\Bigr]
\]
is a $\gamma$-contraction in the sup-norm on the space of bounded value functions.
Value iteration --- the algorithm that repeatedly applies $T$ starting from any
initial $V_0$ --- converges to the unique fixed point $V^*$ (the optimal value
function) precisely because of the Banach fixed point theorem. The convergence
rate $\gamma^n/(1-\gamma)$ gives the iteration complexity of value iteration.
\end{keyresult}

%% ============================================================
%% REFERENCES
%% ============================================================
\section*{References}

\begin{enumerate}[label={[\arabic*]}]
    \item W.~Rudin, \emph{Principles of Mathematical Analysis}, 3rd ed.,
    McGraw-Hill, 1976.

    \item S.~Abbott, \emph{Understanding Analysis}, 2nd ed.,
    Springer, 2015.

    \item G.~B.~Folland, \emph{Real Analysis: Modern Techniques and Their
    Applications}, 2nd ed., Wiley, 1999.

    \item R.~G.~Bartle and D.~R.~Sherbert, \emph{Introduction to Real Analysis},
    4th ed., Wiley, 2011.
\end{enumerate}

\end{document}
